%============================
% SECTION 2: DICOM PARSING
%============================
\section{Reading DICOM Data}

\begin{frame}
\frametitle{DICOM File Structure}
DICOM files are binary with tagged structure:

\begin{itemize}
  \item \textbf{Group and Element tags}: (XXXX, XXXX) in hexadecimal
  \item \textbf{Value Representation (VR)}: Data type (UI, CS, LO, US, OW, etc.)
  \item \textbf{Length}: Bytes in the value
  \item \textbf{Value}: Actual data
\end{itemize}

\vspace{0.5cm}

\textbf{Key DICOM Tags for Image Slices:}
\small
\begin{tabular}{|l|l|}
\hline
\textbf{Tag} & \textbf{Description} \\
\hline
(0028,0010) & Rows (height) \\
(0028,0011) & Columns (width) \\
(0028,0100) & Bits Allocated (8 or 16) \\
(0020,0013) & Instance Number \\
(0020,0032) & Image Position Patient (x, y, z) \\
(7FE0,0010) & Pixel Data (binary image) \\
\hline
\end{tabular}
\end{frame}

\begin{frame}[fragile]{DICOM Parsing: Metadata}
\begin{columns}
\begin{column}{0.5\textwidth}
\begin{lstlisting}
// Extract Header Tags
const rows = dataSet.uint16('x00280010');
const cols = dataSet.uint16('x00280011');
const bits = dataSet.uint16('x00280100');
\end{lstlisting}
\end{column}
\begin{column}{0.5\textwidth}
\small
\textbf{Key Tags:}
\begin{itemize}
    \item \texttt{x00280010}: Image Height (Rows)
    \item \texttt{x00280011}: Image Width (Columns)
    \item \texttt{x00280100}: Bits Allocated (usually 16-bit for CT/MRI)
\end{itemize}
\end{column}
\end{columns}
\end{frame}

\begin{frame}[fragile]{DICOM Parsing: Pixel Data}
\begin{lstlisting}
// Locate the Pixel Data Element (7FE0,0010)
const pixelElement = dataSet.elements.x7fe00010;

// Create a view without copying the underlying buffer
const pixelData = new Uint16Array(
  byteArray.buffer, 
  pixelElement.dataOffset, 
  pixelElement.length / 2
);
\end{lstlisting}
\vfill
\begin{alertblock}{Performance Note}
Using \texttt{Uint16Array} on the existing \texttt{buffer} prevents a full memory copy, which is critical when loading 500+ slices.
\end{alertblock}
\end{frame}

\begin{frame}[fragile]{DICOM Parsing: Spatial Context}
\begin{lstlisting}
// Image Position (Patient) Tag
const posStr = dataSet.string('x00200032');
const [x, y, z] = posStr.split(' ').map(Number);
\end{lstlisting}
\vfill
\begin{itemize}
    \item \textbf{Image Position (Patient):} Specifies the $x, y, z$ coordinates of the upper left hand corner of the slice.
    \item \textbf{Stacking:} We use the \texttt{z} value to sort slices into a coherent 3D volume.
\end{itemize}
\end{frame}

\begin{frame}
\frametitle{Data Type Handling}
\begin{itemize}
  \item \textbf{8-bit images}: Uint8Array (common in radiography)
  \item \textbf{16-bit images}: Uint16Array (common in CT/MRI)
  \item \textbf{Byte ordering}: Little-endian on most systems
  \item \textbf{Pixel representation}: Usually unsigned (0 to max)
  \item Hounsfield Units in CT: Requires proper windowing
\end{itemize}

\vspace{0.5cm}

\textbf{Slice Collection:}
\begin{itemize}
  \item Multiple DICOM files $\rightarrow$ one file per slice
  \item Files must be sorted by spatial position
  \item Use \texttt{ImagePositionPatient} (most reliable)
  \item Fallback to \texttt{InstanceNumber}
\end{itemize}
\end{frame}